% Options for packages loaded elsewhere
\PassOptionsToPackage{unicode}{hyperref}
\PassOptionsToPackage{hyphens}{url}
\documentclass[
  indonesian,
]{article}
\usepackage{xcolor}
\usepackage{amsmath,amssymb}
\setcounter{secnumdepth}{-\maxdimen} % remove section numbering
\usepackage{iftex}
\ifPDFTeX
  \usepackage[T1]{fontenc}
  \usepackage[utf8]{inputenc}
  \usepackage{textcomp} % provide euro and other symbols
\else % if luatex or xetex
  \usepackage{unicode-math} % this also loads fontspec
  \defaultfontfeatures{Scale=MatchLowercase}
  \defaultfontfeatures[\rmfamily]{Ligatures=TeX,Scale=1}
\fi
\usepackage{lmodern}
\ifPDFTeX\else
  % xetex/luatex font selection
\fi
% Use upquote if available, for straight quotes in verbatim environments
\IfFileExists{upquote.sty}{\usepackage{upquote}}{}
\IfFileExists{microtype.sty}{% use microtype if available
  \usepackage[]{microtype}
  \UseMicrotypeSet[protrusion]{basicmath} % disable protrusion for tt fonts
}{}
\makeatletter
\@ifundefined{KOMAClassName}{% if non-KOMA class
  \IfFileExists{parskip.sty}{%
    \usepackage{parskip}
  }{% else
    \setlength{\parindent}{0pt}
    \setlength{\parskip}{6pt plus 2pt minus 1pt}}
}{% if KOMA class
  \KOMAoptions{parskip=half}}
\makeatother
\usepackage{longtable,booktabs,array}
\usepackage{tabularx}
\newcounter{none} % for unnumbered tables
\usepackage{calc} % for calculating minipage widths
% Correct order of tables after \paragraph or \subparagraph
\usepackage{etoolbox}
\makeatletter
\patchcmd\longtable{\par}{\if@noskipsec\mbox{}\fi\par}{}{}
\makeatother
% Allow footnotes in longtable head/foot
\IfFileExists{footnotehyper.sty}{\usepackage{footnotehyper}}{\usepackage{footnote}}
\makesavenoteenv{longtable}
\ifLuaTeX
\usepackage[bidi=basic,shorthands=off,provide=*]{babel}
\else
\usepackage[bidi=default,shorthands=off,provide=*]{babel}
\fi
\ifLuaTeX
  \usepackage{selnolig} % disable illegal ligatures
\fi
\setlength{\emergencystretch}{3em} % prevent overfull lines
\providecommand{\tightlist}{%
  \setlength{\itemsep}{0pt}\setlength{\parskip}{0pt}}
\usepackage{bookmark}
\usepackage{imakeidx}
\makeindex[title=Indeks istilah dan situasi,intoc,columns=2]
\IfFileExists{xurl.sty}{\usepackage{xurl}}{} % add URL line breaks if available
\urlstyle{same}
% fallback for those not using the hyperref driver hyperxmp:
\makeatletter
\@ifundefined{xmpquote}{\newcommand{\xmpquote}[1]{#1}}{}
\makeatother
\hypersetup{
  pdftitle={Panduan Komunikasi TGA dan Skripsi: Retno Widodo Dwi Pramono},
  pdflang={id-ID},
  hidelinks,
  pdfcreator={LaTeX via pandoc}}

\usepackage[a4paper,top=22mm,bottom=22mm,left=23mm,right=20mm,headheight=15pt,headsep=7mm,footskip=11mm]{geometry}
\usepackage{setspace}
\usepackage{titlesec}
\usepackage{enumitem}
\usepackage{fancyhdr}
\usepackage{needspace}
\usepackage{ragged2e}
\usepackage[most]{tcolorbox}
\usepackage{fvextra}
\RecustomVerbatimEnvironment{verbatim}{Verbatim}{%
  breaklines=true,
  breakanywhere=true,
  fontsize=\small,
  xleftmargin=1em,
  xrightmargin=1em}

\definecolor{Ink}{HTML}{252525}
\definecolor{Accent}{HTML}{355C7D}
\definecolor{Soft}{HTML}{F4F7F9}
\definecolor{Rule}{HTML}{D5DEE5}
\definecolor{LinkInk}{HTML}{274C61}

\setstretch{1.08}
\setlength{\parindent}{0pt}
\setlength{\parskip}{0.42em}
\setlength{\emergencystretch}{3em}
\raggedbottom
\clubpenalty=10000
\widowpenalty=10000
\displaywidowpenalty=10000

\setlist[itemize]{leftmargin=1.65em,itemsep=0.2em,topsep=0.3em,parsep=0pt}
\setlist[enumerate]{leftmargin=1.85em,itemsep=0.2em,topsep=0.3em,parsep=0pt}

\titleformat{\section}
  {\normalfont\Large\bfseries\color{Ink}}
  {}{0pt}{}
\titlespacing*{\section}{0pt}{1.35em}{0.5em}
\titleformat{\subsection}
  {\normalfont\large\bfseries\color{Ink}}
  {}{0pt}{}
\titlespacing*{\subsection}{0pt}{1.0em}{0.35em}
\titleformat{\subsubsection}
  {\normalfont\normalsize\bfseries\color{Ink}}
  {}{0pt}{}
\titlespacing*{\subsubsection}{0pt}{0.85em}{0.25em}

% The Markdown document uses level-2 and level-3 headings for its main
% sections. Promote those levels by one step in the printable edition so
% the table of contents has a useful hierarchy.
\let\OriginalSubsection\subsection
\let\OriginalSubsubsection\subsubsection
\renewcommand{\subsection}{\section}
\renewcommand{\subsubsection}{\OriginalSubsection}

\pagestyle{fancy}
\fancyhf{}
\fancyhead[L]{\footnotesize\textit{Panduan komunikasi TGA dan skripsi}}
\fancyhead[R]{\footnotesize\textit{Retno Widodo Dwi Pramono}}
\fancyfoot[C]{\thepage}
\renewcommand{\headrulewidth}{0.3pt}
\renewcommand{\footrulewidth}{0pt}
\fancypagestyle{plain}{%
  \fancyhf{}
  \fancyfoot[C]{\thepage}
  \renewcommand{\headrulewidth}{0pt}
}

\AtBeginEnvironment{quote}{%
  \small
  \setlength{\leftskip}{1.25em}
  \setlength{\rightskip}{1.25em}
  \color{Ink}%
}
\AtBeginEnvironment{longtable}{%
  \footnotesize
  \setlength{\tabcolsep}{4pt}
  \renewcommand{\arraystretch}{1.16}
  \sloppy
}
\setlength{\LTpre}{0.55em}
\setlength{\LTpost}{0.8em}
\renewcommand{\arraystretch}{1.16}

\newcommand{\quicklink}[2]{\hyperref[#1]{\textcolor{LinkInk}{#2}}}
\newcommand{\keywordterm}[1]{\textbf{#1}\index{#1}}

\title{Panduan Komunikasi TGA dan Skripsi: Retno Widodo Dwi Pramono}
\author{}
\date{18 September 2026}


\providecommand{\passthrough}[1]{#1}

\providecommand{\passthrough}[1]{#1}
\begin{document}
\maketitle
\renewcommand*\contentsname{Daftar Isi}
\setcounter{tocdepth}{2}
\tableofcontents\clearpage
\begin{quote}
\textbf{Status:} panduan kerja dan kumpulan templat komunikasi. Templat
tidak perlu disalin kata demi kata. Sesuaikan isinya dengan hubungan,
kanal, dan kebiasaan penerima.

\textbf{Dasar penyesuaian:} Pemetaan Dosen dan Riset-Compass PWK,
Roadmap Skripsi TGA - Borrowed Size dan Agglomeration Shadow Jakarta,
serta kebutuhan komunikasi administratif mahasiswa.

\textbf{Aturan umum:} tidak menggunakan pembuka ``Assalamualaikum''.
Setiap kelompok memiliki contoh untuk percakapan langsung dan WhatsApp.
Contoh email hanya tersedia untuk urusan administrasi akademik DTAP.
\end{quote}

\section{1. Cara menggunakan panduan}\label{cara-menggunakan-panduan}

\begin{itemize}
\tightlist
\item
  Ganti semua bagian dalam kurung siku, misalnya \texttt{{[}Nama{]}},
  \texttt{{[}NIM{]}}, \texttt{{[}tanggal{]}}, dan \texttt{{[}tautan{]}}.
\item
  Versi ``bicara langsung'' berisi kerangka singkat, bukan naskah yang
  harus dibaca seperti surat.
\item
  Untuk WhatsApp, sesuaikan panjang pesan dengan konteks dan usahakan
  mengirim satu pesan yang utuh.
\item
  Gunakan ``Bapak Retno'' ketika menyebut beliau dalam kalimat. Untuk
  menyapa beliau secara langsung, ``Pak Retno'' terdengar lebih wajar
  dan tetap sopan.
\item
  Email hanya digunakan untuk urusan administrasi akademik DTAP.
\item
  Gunakan profil riset sebagai dugaan awal tentang hal-hal yang mungkin
  beliau anggap penting dalam proposal. Jangan menganggapnya sebagai
  gambaran pasti tentang kepribadian, jadwal, atau kecepatan respons
  beliau.
\end{itemize}

\section{2. Pegangan awal untuk berkomunikasi dengan Bapak
Retno}\label{pegangan-awal-untuk-berkomunikasi-dengan-bapak-retno}

Dari pemetaan dosen, terlihat beberapa tema yang sering muncul dalam
riset Bapak Retno: evaluasi pembangunan, hubungan desa-kota, kebijakan
pertanahan, keadilan, kesejahteraan, dan metode yang pragmatis. Namun,
gambaran tentang gaya bimbingan dan keterbatasan waktu beliau masih
berupa dugaan. Karena itu, gunakan beberapa pegangan berikut sebagai
titik awal:

\begin{itemize}
\tightlist
\item
  Tunjukkan hubungan yang jelas antara masalah, teori, metode, data, dan
  hasil yang dituju.
\item
  Tunjukkan keputusan yang sudah dibuat dan bagian yang masih memerlukan
  arahan.
\item
  Sampaikan keterbatasan data sejak awal. Jangan menyembunyikan risiko
  hanya agar proposal terlihat lebih mudah dikerjakan.
\item
  Bahas dampak, pembagian manfaat dan beban, atau implikasi perencanaan
  hanya jika didukung oleh bukti.
\item
  Jangan mengganti persiapan riset dengan pujian terhadap publikasi
  beliau.
\item
  Jangan menganggap beliau pasti menerima, membalas cepat, menyukai
  metode tertentu, atau bersedia mengubah jadwal tanpa konfirmasi
  langsung.
\end{itemize}

\section{3. Matriks nada dan kanal}\label{matriks-nada-dan-kanal}

{\def\LTcaptype{none} % do not increment counter
\begin{longtable}[]{@{}
  >{\raggedright\arraybackslash}p{(\linewidth - 6\tabcolsep) * \real{0.2500}}
  >{\raggedright\arraybackslash}p{(\linewidth - 6\tabcolsep) * \real{0.2500}}
  >{\raggedright\arraybackslash}p{(\linewidth - 6\tabcolsep) * \real{0.2500}}
  >{\raggedright\arraybackslash}p{(\linewidth - 6\tabcolsep) * \real{0.2500}}@{}}
\toprule\noalign{}
\begin{minipage}[b]{\linewidth}\raggedright
Penerima
\end{minipage} & \begin{minipage}[b]{\linewidth}\raggedright
Komunikasi langsung
\end{minipage} & \begin{minipage}[b]{\linewidth}\raggedright
WhatsApp
\end{minipage} & \begin{minipage}[b]{\linewidth}\raggedright
Email
\end{minipage} \\
\midrule\noalign{}
\endhead
\bottomrule\noalign{}
\endlastfoot
Bapak Retno, dosen pembimbing & Hormat, tenang, dan terstruktur tanpa
merendahkan diri & Ringkas, sopan, langsung pada agenda dan permintaan &
Tidak digunakan \\
Administrasi akademik DTAP & Formal dan faktual & Formal, disertai
identitas dan pertanyaan yang jelas & Formal, dengan subjek dan lampiran
yang jelas \\
Teman seangkatan, angkatan 23, dan pra-TA angkatan 24 & Sopan, santai,
tidak menggurui & Ramah, jelas, dan menghargai waktu & Tidak
digunakan \\
\end{longtable}
}

\subsection{3.1 Pola dasar komunikasi
langsung}\label{pola-dasar-komunikasi-langsung}

\begin{enumerate}
\def\labelenumi{\arabic{enumi}.}
\tightlist
\item
  Sapa sesuai waktu.
\item
  Sebutkan identitas jika belum saling mengenal atau jika situasinya
  formal.
\item
  Jelaskan konteks dalam satu atau dua kalimat.
\item
  Sampaikan satu permintaan atau keputusan yang ingin dibahas.
\item
  Tutup dengan terima kasih dan tindak lanjut.
\end{enumerate}

Contoh umum kepada dosen:

\begin{quote}
Selamat pagi, Pak. Perkenalkan, saya {[}Nama{]}, mahasiswa {[}program
studi{]} dengan NIM {[}NIM{]}. Saya ingin meminta arahan tentang {[}satu
keperluan{]}. Saya sudah menyiapkan {[}bahan atau hasil{]}. Pertanyaan
saya: {[}pertanyaan atau pilihan{]}. Terima kasih, Pak.
\end{quote}

\subsection{3.2 Pola dasar WhatsApp}\label{pola-dasar-whatsapp}

\begin{verbatim}
Selamat pagi, Pak Retno. Perkenalkan, saya [Nama] ([NIM]).

[Konteks singkat]. Saya ingin meminta arahan tentang [keperluan]. Saya sudah menyiapkan [berkas atau hasil]. Apakah Bapak bisa [permintaan atau pilihan waktu]?

Terima kasih, Pak.
\end{verbatim}

\section{4. Siklus komunikasi dengan dosen
pembimbing}\label{siklus-komunikasi-dengan-dosen-pembimbing}

{\def\LTcaptype{none} % do not increment counter
\begin{longtable}[]{@{}
  >{\raggedright\arraybackslash}p{(\linewidth - 4\tabcolsep) * \real{0.3333}}
  >{\raggedright\arraybackslash}p{(\linewidth - 4\tabcolsep) * \real{0.3333}}
  >{\raggedright\arraybackslash}p{(\linewidth - 4\tabcolsep) * \real{0.3333}}@{}}
\toprule\noalign{}
\begin{minipage}[b]{\linewidth}\raggedright
Fase
\end{minipage} & \begin{minipage}[b]{\linewidth}\raggedright
Komunikasi yang mungkin terjadi
\end{minipage} & \begin{minipage}[b]{\linewidth}\raggedright
Tujuan
\end{minipage} \\
\midrule\noalign{}
\endhead
\bottomrule\noalign{}
\endlastfoot
Pra-TA & Pendekatan awal, pengiriman ringkasan, tindak lanjut & Meminta
waktu untuk berdiskusi dan menanyakan kemungkinan bimbingan \\
Awal bimbingan & Perkenalan, penetapan fokus, penjadwalan & Menyepakati
cara kerja dan keputusan awal \\
Proposal & Pengiriman draf, permintaan umpan balik, klarifikasi &
Menghasilkan proposal yang koheren \\
Sumber terbuka dan proksi & Membahas konstruksi indikator, hasil
akuisisi, dan validasi & Menjaga pengukuran dapat ditelusuri dan
direplikasi \\
Masa tunggu & Mengerjakan baseline terbuka, meninjau literatur, atau
menyiapkan keputusan metode & Tetap bekerja tanpa mengklaim data yang
belum tersedia \\
Analisis & Membawa hasil sementara, kendala, dan sensitivitas & Menguji
bukti dan batas klaim \\
Progres dan prasidang & Meminta persetujuan, memeriksa catatan
bimbingan, dan menyiapkan berkas & Memenuhi persyaratan akademik dan
administratif \\
Akhir & Mengabarkan kelulusan, menyelesaikan revisi, dan berterima kasih
& Menutup hubungan kerja dengan baik \\
\end{longtable}
}

\subsection{4.1 Pendekatan awal sebelum menjadi
pembimbing}\label{pendekatan-awal-sebelum-menjadi-pembimbing}

\textbf{Bicara langsung}

\begin{quote}
Selamat pagi, Pak Retno. Perkenalkan, saya {[}Nama{]}, mahasiswa
{[}program studi{]} angkatan 2022. Saya sedang menyiapkan pra-TA tentang
{[}topik singkat{]}. Saya sudah membuat ringkasan satu halaman yang
berisi masalah, pertanyaan, metode awal, dan kebutuhan data. Jika Bapak
berkenan, saya ingin meminta waktu untuk mendiskusikan topik ini
sekaligus menanyakan kemungkinan bimbingan.
\end{quote}

\textbf{WhatsApp}

\begin{verbatim}
Selamat pagi, Pak Retno. Perkenalkan, saya [Nama] (NIM [NIM]), mahasiswa PWK angkatan 2022.

Saya sedang menyiapkan pra-TA tentang [topik singkat] dan sudah membuat ringkasan satu halaman yang berisi masalah, pertanyaan, metode awal, serta kebutuhan data. Apakah saya boleh mengirimkannya dan meminta waktu untuk berdiskusi tentang kemungkinan bimbingan?

Terima kasih, Pak.
\end{verbatim}

\subsection{4.2 Mengirim ringkasan setelah mendapat
izin}\label{mengirim-ringkasan-setelah-mendapat-izin}

\textbf{Bicara langsung}

\begin{quote}
Terima kasih, Pak. Ini ringkasan yang tadi saya sampaikan. Saya terutama
ingin berkonsultasi tentang kelayakan pertanyaan, batas penelitian, dan
kecukupan metode. Bagian yang masih berupa asumsi sudah saya tandai.
\end{quote}

\textbf{WhatsApp}

\begin{verbatim}
Sesuai izin Bapak, saya kirim ringkasan pra-TA saya: [nama berkas atau tautan].

Bagian yang terutama ingin saya konsultasikan adalah kelayakan pertanyaan, batas penelitian, dan kecukupan metode. Asumsi serta risiko datanya sudah saya tandai di dalam berkas.
\end{verbatim}

\subsection{4.3 Menindaklanjuti pendekatan tanpa
menekan}\label{menindaklanjuti-pendekatan-tanpa-menekan}

\textbf{Bicara langsung}

\begin{quote}
Pak, saya ingin menindaklanjuti ringkasan yang saya kirim pada
{[}tanggal{]}. Jika Bapak bersedia mendiskusikannya, saya siap
menyesuaikan dengan waktu yang tersedia.
\end{quote}

\textbf{WhatsApp}

\begin{verbatim}
Selamat pagi, Pak Retno. Saya ingin menindaklanjuti ringkasan pra-TA yang saya kirim pada [tanggal]. Jika Bapak bersedia mendiskusikannya, saya siap menyesuaikan dengan waktu yang tersedia.

Terima kasih, Pak.
\end{verbatim}

\subsection{4.4 Jika Bapak Retno bersedia
membimbing}\label{jika-bapak-retno-bersedia-membimbing}

\textbf{Bicara langsung}

\begin{quote}
Terima kasih sudah bersedia membimbing saya, Pak. Pada pertemuan awal,
saya ingin menyepakati target kerja, bentuk draf, kanal komunikasi, dan
cara menjadwalkan bimbingan agar alurnya jelas.
\end{quote}

\textbf{WhatsApp}

\begin{verbatim}
Terima kasih sudah bersedia membimbing saya, Pak.

Pada pertemuan awal nanti, saya ingin menyepakati target kerja, bentuk draf, kanal komunikasi, dan cara menjadwalkan bimbingan agar alurnya jelas.
\end{verbatim}

\subsection{4.5 Jika beliau belum dapat atau tidak bersedia
membimbing}\label{jika-beliau-belum-dapat-atau-tidak-bersedia-membimbing}

\textbf{Bicara langsung}

\begin{quote}
Baik, Pak. Terima kasih sudah memberi kabar dan meluangkan waktu. Saya
memahami keputusan Bapak. Saya akan mengikuti prosedur yang berlaku
untuk mencari pembimbing lain.
\end{quote}

\textbf{WhatsApp}

\begin{verbatim}
Baik, Pak. Terima kasih sudah memberi kabar dan meluangkan waktu. Saya memahami keputusan Bapak dan akan mengikuti prosedur yang berlaku untuk mencari pembimbing lain.
\end{verbatim}

\subsection{4.6 Mengusulkan jadwal
bimbingan}\label{mengusulkan-jadwal-bimbingan}

\textbf{Bicara langsung}

\begin{quote}
Pak, untuk bimbingan berikutnya, saya punya tiga usulan waktu: {[}hari
dan tanggal{]} pukul {[}jam 1{]}, {[}hari dan tanggal{]} pukul {[}jam
2{]}, atau {[}hari dan tanggal{]} pukul {[}jam 3{]}. Bimbingannya
kira-kira berlangsung selama {[}durasi{]} dengan agenda {[}agenda{]}.
Apakah salah satunya sesuai bagi Bapak?
\end{quote}

\textbf{WhatsApp}

\begin{verbatim}
Selamat pagi, Pak Retno. Untuk bimbingan berikutnya, saya ingin mengusulkan beberapa pilihan waktu:

1. [Hari, tanggal], pukul [jam]
2. [Hari, tanggal], pukul [jam]
3. [Hari, tanggal], pukul [jam]

Bimbingannya kira-kira berlangsung selama [durasi] dengan agenda [agenda]. Saya akan membawa [berkas atau bahan]. Apakah salah satu waktu tersebut sesuai bagi Bapak? Jika belum, saya siap menyesuaikan.
\end{verbatim}

\subsection{4.7 Mengonfirmasi jadwal yang sudah
disepakati}\label{mengonfirmasi-jadwal-yang-sudah-disepakati}

\textbf{Bicara langsung}

\begin{quote}
Baik, Pak. Saya catat bimbingannya pada {[}hari, tanggal{]}, pukul
{[}jam{]}, di {[}tempat atau tautan{]}, dengan agenda {[}agenda{]}. Saya
akan mengirim bahannya paling lambat {[}waktu{]}.
\end{quote}

\textbf{WhatsApp}

\begin{verbatim}
Baik, Pak. Saya konfirmasi bimbingannya pada [hari, tanggal], pukul [jam], di [tempat atau tautan]. Agenda yang saya catat adalah [agenda]. Bahannya akan saya kirim paling lambat [waktu].
\end{verbatim}

\subsection{4.8 Jika dosen mengubah jadwal yang sudah
disepakati}\label{jika-dosen-mengubah-jadwal-yang-sudah-disepakati}

\textbf{Bicara langsung}

\begin{quote}
Baik, Pak. Untuk jadwal penggantinya, apakah {[}opsi 1{]} atau {[}opsi
2{]} memungkinkan? Jika belum, saya bisa menyesuaikan dengan waktu lain.
\end{quote}

\textbf{WhatsApp}

\begin{verbatim}
Baik, Pak. Untuk jadwal penggantinya, apakah salah satu waktu berikut memungkinkan?

1. [Hari, tanggal], pukul [jam]
2. [Hari, tanggal], pukul [jam]

Jika belum sesuai, saya siap menyesuaikan waktu lain.
\end{verbatim}

Jangan menulis ``Bukankah sebelumnya sudah disepakati?'' atau
menyiratkan bahwa perubahan jadwal merupakan kesalahan pribadi. Cukup
pastikan jadwal baru dan catat perubahannya.

\subsection{4.9 Jika mahasiswa yang perlu mengubah
jadwal}\label{jika-mahasiswa-yang-perlu-mengubah-jadwal}

\textbf{Bicara langsung}

\begin{quote}
Mohon maaf, Pak. Saya perlu menjadwalkan ulang bimbingan karena
{[}alasan singkat dan wajar{]}. Sebagai pengganti, saya mengusulkan
{[}opsi 1{]} atau {[}opsi 2{]}. Apakah salah satunya sesuai bagi Bapak?
\end{quote}

\textbf{WhatsApp}

\begin{verbatim}
Selamat pagi, Pak Retno. Mohon maaf, saya perlu menjadwalkan ulang bimbingan pada [hari, tanggal] karena [alasan singkat].

Sebagai pengganti, saya mengusulkan dua pilihan waktu:
1. [Hari, tanggal], pukul [jam]
2. [Hari, tanggal], pukul [jam]

Apakah salah satunya sesuai bagi Bapak? Jika belum, saya siap menyesuaikan. Terima kasih, Pak.
\end{verbatim}

\subsection{4.10 Mengirim draf sebelum
bimbingan}\label{mengirim-draf-sebelum-bimbingan}

\textbf{Bicara langsung}

\begin{quote}
Pak, ini draf versi {[}nomor{]}, tanggal {[}tanggal{]}. Saya terutama
memerlukan arahan pada {[}bagian 1{]} dan {[}bagian 2{]}. Perubahan dari
versi sebelumnya sudah saya tandai. Hal-hal yang perlu diputuskan saya
letakkan di halaman {[}nomor{]}.
\end{quote}

\textbf{WhatsApp}

\begin{verbatim}
Selamat sore, Pak Retno. Saya kirim draf [nama berkas], versi [nomor], tanggal [tanggal]: [tautan].

Saya terutama memerlukan arahan tentang hal-hal berikut:
1. [bagian atau keputusan 1]
2. [bagian atau keputusan 2]

Perubahan dari versi sebelumnya sudah saya tandai. Terima kasih, Pak.
\end{verbatim}

\subsection{4.11 Meminta umpan balik tanpa menuntut balasan
segera}\label{meminta-umpan-balik-tanpa-menuntut-balasan-segera}

\textbf{Bicara langsung}

\begin{quote}
Pak, jika Bapak sudah sempat membaca drafnya, saya ingin meminta arahan
tentang {[}satu bagian{]}. Jika belum, saya akan menyesuaikan agenda
bimbingan berikutnya.
\end{quote}

\textbf{WhatsApp}

\begin{verbatim}
Selamat pagi, Pak Retno. Saya ingin menindaklanjuti draf yang saya kirim pada [tanggal]. Jika Bapak sudah sempat membacanya, saya terutama memerlukan arahan tentang [bagian atau keputusan]. Jika belum, saya akan menyesuaikan agenda bimbingan berikutnya.
\end{verbatim}

\subsection{4.12 Mengklarifikasi arahan yang belum jelas atau tampak
bertentangan}\label{mengklarifikasi-arahan-yang-belum-jelas-atau-tampak-bertentangan}

\textbf{Bicara langsung}

\begin{quote}
Pak, saya ingin memastikan bahwa saya tidak salah menerapkan arahan.
Pada bimbingan tanggal {[}tanggal{]}, saya mencatat {[}arahan A{]},
tetapi di bagian lain saya mencatat {[}arahan B{]}. Untuk draf
berikutnya, apakah saya sebaiknya memilih {[}pilihan 1{]} atau
{[}pilihan 2{]}?
\end{quote}

\textbf{WhatsApp}

\begin{verbatim}
Pak, saya ingin memastikan bahwa saya tidak salah menerapkan arahan.

Pada bimbingan tanggal [tanggal], saya mencatat [arahan A]. Namun, di bagian lain saya mencatat [arahan B]. Untuk draf berikutnya, apakah saya sebaiknya memilih [pilihan 1] atau [pilihan 2]?

Terima kasih atas arahannya, Pak.
\end{verbatim}

Jangan menulis ``arahan Bapak berubah-ubah''. Tunjukkan catatan yang
Anda miliki, lalu tanyakan pilihan mana yang perlu dipakai.

\subsection{4.13 Membahas arah data dan
metode}\label{membahas-arah-data-dan-metode}

\textbf{Bicara langsung}

\begin{quote}
Pak, saya menggunakan proksi dan sumber terbuka sebagai fondasi utama,
termasuk data pemerintah yang terbuka apabila sesuai. Desainnya
multitemporal sesuai bukti. Saya ingin membahas konstruksi indikator,
kesebandingan tahun, dan validasinya.
\end{quote}

\textbf{WhatsApp}

\begin{verbatim}
Pak Retno, saya sudah memetakan kebutuhan pengukuran lokasi usaha, atribut fasilitas, jaringan, informasi operator, dan massa lokal. Rancangan tidak bergantung pada permohonan data tertutup.

Saya ingin membahas prioritas sumber yang benar-benar tersedia dan asumsi untuk estimasi pekerjaan atau relasi metropolitan. Ringkasan sumber, periode, dan keterbatasannya: [tautan].
\end{verbatim}

\subsection{4.14 Melaporkan hasil akuisisi dan
audit}\label{melaporkan-hasil-akuisisi-dan-audit}

\begin{verbatim}
Pak Retno, berikut hasil akuisisi sumber terbuka:

- [Sumber]: [berkas, unit, periode, cakupan, dan mutu yang telah diperiksa].
- [Sumber]: [hasil dan kendala].

Saya mengusulkan [konstruksi indikator] dengan asumsi [asumsi] dan pemeriksaan independen [pemeriksaan]. Bagian yang belum dapat diverifikasi adalah [bagian]. Saya ingin membahas dampaknya bagi analisis.
\end{verbatim}

\subsection{4.15 Melaporkan progres konstruksi
data}\label{melaporkan-progres-konstruksi-data}

\begin{verbatim}
Pak Retno, saya telah menyelesaikan [ekstraksi, pencocokan lokasi, audit seri, atau estimasi awal]. Hasilnya tersedia di [tautan].

Pada bimbingan berikutnya, saya ingin membahas [isu metode] dan pilihan [unit, koefisien, rentang tahun, atau validasi]. Modul yang sudah layak tetap berjalan, sementara keterbatasan modul lain dicatat secara terbuka.
\end{verbatim}

\subsection{4.16 Laporan kemajuan
singkat}\label{laporan-kemajuan-singkat}

\textbf{Bicara langsung}

\begin{quote}
Pak, sejak bimbingan terakhir, saya sudah menyelesaikan {[}pekerjaan{]}.
Kendala utamanya adalah {[}kendala{]}. Berikutnya, saya akan
{[}langkah{]}. Saya memerlukan arahan Bapak tentang {[}keputusan{]}.
\end{quote}

\textbf{WhatsApp}

\begin{verbatim}
Pembaruan singkat, Pak:

Sudah selesai: [pekerjaan atau hasil].
Kendala: [kendala yang spesifik].
Langkah berikutnya: [langkah].
Arahan yang saya perlukan: [satu keputusan].
\end{verbatim}

\subsection{4.17 Jika terlambat, sakit, atau tidak dapat
hadir}\label{jika-terlambat-sakit-atau-tidak-dapat-hadir}

\textbf{Bicara langsung}

\begin{quote}
Mohon maaf, Pak. Saya tidak dapat hadir pada {[}waktu{]} karena
{[}alasan singkat{]}. Sebagai pengganti, saya ingin mengusulkan {[}opsi
1{]} atau {[}opsi 2{]}. Apakah salah satunya sesuai bagi Bapak?
\end{quote}

\textbf{WhatsApp}

\begin{verbatim}
Selamat pagi, Pak Retno. Mohon maaf, saya tidak dapat mengikuti bimbingan pada [hari, tanggal] karena [alasan singkat].

Sebagai pengganti, saya ingin mengusulkan [opsi 1] atau [opsi 2]. Apakah salah satunya sesuai bagi Bapak? Saya tetap akan mengirim [bahan atau hasil] agar pekerjaan tetap berjalan.

Terima kasih, Pak.
\end{verbatim}

\textbf{Jika terlambat}

\textbf{Bicara langsung}

\begin{quote}
Mohon maaf, Pak, saya terlambat sekitar {[}durasi{]} karena {[}alasan
singkat{]}. Apakah bimbingannya masih dapat dilanjutkan? Jika tidak,
saya siap menjadwalkan ulang.
\end{quote}

\textbf{WhatsApp}

\begin{verbatim}
Selamat pagi, Pak Retno. Mohon maaf, saya kemungkinan terlambat sekitar [durasi] karena [alasan singkat]. Apakah bimbingannya masih dapat dilanjutkan? Jika tidak, saya siap menjadwalkan ulang.
\end{verbatim}

\subsection{4.18 Progres 1 dan Progres 2}\label{progres-1-dan-progres-2}

\textbf{Bicara langsung}

\begin{quote}
Pak, saya sedang menyiapkan Progres {[}1 atau 2{]}. Bahan yang sudah
tersedia adalah {[}daftar singkat{]}, sedangkan bagian yang belum
selesai adalah {[}bagian{]}. Apa yang perlu saya prioritaskan sebelum
Progres tersebut?
\end{quote}

\textbf{WhatsApp}

\begin{verbatim}
Pak Retno, saya sedang menyiapkan Progres [1 atau 2].

Bahan yang sudah tersedia: [bahan atau hasil].
Bagian yang belum selesai: [bagian].
Apa yang perlu saya prioritaskan sebelum Progres tersebut?

Terima kasih, Pak.
\end{verbatim}

\subsection{4.19 Meminta persetujuan menuju
prasidang}\label{meminta-persetujuan-menuju-prasidang}

\textbf{Bicara langsung}

\begin{quote}
Pak, saya ingin memastikan bahwa naskah dan berkas saya sudah siap
sebelum mendaftar prasidang. Saya sudah menyiapkan naskah, PDF, hasil
Turnitin, nilai Progres, dan catatan bimbingan. Saat ini, jumlah
bimbingan yang tercatat adalah {[}jumlah{]}. Apakah masih ada perbaikan
isi atau administrasi yang perlu saya selesaikan?
\end{quote}

\textbf{WhatsApp}

\begin{verbatim}
Pak Retno, saya ingin memastikan bahwa naskah dan berkas saya sudah siap sebelum mendaftar prasidang.

Saya sudah menyiapkan naskah, PDF, hasil Turnitin, nilai Progres, dan catatan bimbingan. Saat ini, jumlah bimbingan yang tercatat adalah [jumlah].

Apakah masih ada perbaikan isi atau administrasi yang perlu saya selesaikan sebelum mendaftar?

Terima kasih, Pak.
\end{verbatim}

\subsection{4.20 Ucapan terima kasih setelah tahap
akhir}\label{ucapan-terima-kasih-setelah-tahap-akhir}

\textbf{Bicara langsung}

\begin{quote}
Pak, terima kasih atas waktu, arahan, dan kritik Bapak selama proses
TGA. Masukan Bapak membantu saya memperbaiki {[}aspek konkret{]}. Saya
akan menyelesaikan tindak lanjutnya dan memastikan klaim akhir tidak
melampaui bukti yang tersedia.
\end{quote}

\textbf{WhatsApp}

\begin{verbatim}
Pak Retno, terima kasih atas waktu, arahan, dan kritik Bapak selama proses TGA. Masukan Bapak sangat membantu saya memperbaiki [aspek konkret].

Saya akan menyelesaikan tindak lanjutnya dan memastikan klaim akhir tidak melampaui bukti yang tersedia. Terima kasih banyak, Pak.
\end{verbatim}

\section{5. Komunikasi dengan administrasi akademik
DTAP}\label{komunikasi-dengan-administrasi-akademik-dtap}

\subsection{5.1 Prinsip khusus}\label{prinsip-khusus}

\begin{itemize}
\tightlist
\item
  Gunakan nama lengkap, NIM, program studi, dan konteks TGA pada pesan
  pertama.
\item
  Berikan konteks secukupnya, lalu tanyakan prosedur, dokumen, kanal,
  atau tenggat yang belum jelas.
\item
  Bedakan aturan resmi dari informasi yang didengar dari kakak tingkat.
\item
  Gunakan kanal resmi untuk pendaftaran, surat, nilai, Turnitin,
  prasidang, sidang, dan yudisium.
\item
  Periksa kanal rujukan TGA PWK UGM di
  \href{https://sites.google.com/ugm.ac.id/tgapwkugm}{laman resmi TGA}
  dan gunakan informasi dengan tanggal pembaruan paling baru.
\item
  Sebutkan setiap lampiran dengan nama berkas, versi, dan tanggal yang
  jelas.
\item
  Jika belum ada jawaban, ikuti waktu layanan yang tercantum. Jika tidak
  ada, beri jeda yang wajar sebelum mengirim satu pesan tindak lanjut.
\item
  Jangan mengirim data pribadi mahasiswa lain atau dokumen internal
  tanpa izin.
\end{itemize}

\subsection{5.2 Templat umum untuk percakapan
langsung}\label{templat-umum-untuk-percakapan-langsung}

\begin{quote}
Selamat pagi, Bapak/Ibu. Perkenalkan, saya {[}Nama{]}, NIM {[}NIM{]},
mahasiswa {[}program studi{]} angkatan 2022. Saya ingin menanyakan
{[}satu urusan{]}. Saya sudah menyiapkan {[}dokumen pendukung atau nomor
permohonan{]}. Mohon informasi tentang langkah berikutnya dan kanal yang
harus saya gunakan. Terima kasih.
\end{quote}

\subsection{5.3 Templat umum untuk
WhatsApp}\label{templat-umum-untuk-whatsapp}

\begin{verbatim}
Selamat pagi, Bapak/Ibu. Perkenalkan, saya [Nama], NIM [NIM], mahasiswa [program studi] angkatan 2022.

Saya ingin menanyakan [satu urusan]. Saat ini, [konteks singkat]. Dokumen pendukung atau nomor permohonannya adalah [nama atau nomor].

Apakah saya perlu [langkah yang ingin dikonfirmasi] atau ada prosedur lain yang harus saya ikuti? Terima kasih.
\end{verbatim}

\subsection{5.4 Templat umum untuk
email}\label{templat-umum-untuk-email}

\textbf{Subjek:}
\texttt{{[}TGA{]}\ Permohonan\ informasi\ {[}urusan{]}:\ {[}Nama{]},\ {[}NIM{]}}

\begin{verbatim}
Yth. Tim Administrasi Akademik DTAP,

Perkenalkan, saya:
Nama          : [Nama lengkap]
NIM           : [NIM]
Program studi : [Program studi]
Angkatan      : 2022
Keperluan     : [urusan]

Saya ingin menanyakan [pertanyaan yang spesifik]. Saat ini, [konteks singkat].

Dokumen pendukung: [nama dokumen]
Nomor permohonan, jika ada: [nomor]

Mohon informasi tentang:
1. [pertanyaan atau prosedur 1]
2. [pertanyaan atau prosedur 2]
3. [dokumen atau kanal yang harus digunakan]

Terima kasih atas bantuannya.

Hormat saya,

[Nama lengkap]
[NIM]
[Nomor telepon jika diperlukan]
\end{verbatim}

\subsection{5.5 Menanyakan kelayakan administratif
TGA}\label{menanyakan-kelayakan-administratif-tga}

\textbf{Langsung}

\begin{quote}
Saya ingin memastikan kelayakan administratif saya untuk mengajukan TGA.
Mohon informasi tentang dokumen dan status yang perlu saya periksa,
terutama terkait ujian komprehensif, mata kuliah, KRS, dan MK Tugas
Akhir.
\end{quote}

\textbf{WhatsApp}

\begin{verbatim}
Mohon informasi, Bapak/Ibu. Saya ingin memeriksa kelayakan administratif untuk mengajukan TGA. Dokumen atau status apa saja yang perlu saya pastikan, terutama terkait ujian komprehensif, mata kuliah semester 1 hingga 7, KRS, dan MK Tugas Akhir 6 SKS?
\end{verbatim}

\textbf{Email}

\begin{verbatim}
Yth. Tim Administrasi Akademik DTAP,

Saya ingin memeriksa kelayakan administratif untuk mengajukan TGA. Mohon informasi tentang persyaratan dan cara memeriksa status ujian komprehensif, mata kuliah yang telah ditempuh, KRS, serta MK Tugas Akhir.

Terima kasih.

Hormat saya,
[Nama]
[NIM]
\end{verbatim}

\subsection{5.6 Menanyakan pendaftaran, jadwal, atau laman
terbaru}\label{menanyakan-pendaftaran-jadwal-atau-laman-terbaru}

\textbf{Langsung}

\begin{quote}
Mohon informasi tentang kanal dan jadwal resmi pendaftaran {[}TGA,
prasidang, atau sidang{]}. Saya menemukan informasi di {[}sumber{]},
yang saya akses pada {[}tanggal{]}. Apakah informasi tersebut masih
berlaku dan apakah ada formulir yang lebih baru?
\end{quote}

\textbf{WhatsApp}

\begin{verbatim}
Mohon informasi tentang kanal dan jadwal resmi pendaftaran [TGA, prasidang, atau sidang]. Saya menemukan informasi di [sumber], yang saya akses pada [tanggal]. Apakah informasi tersebut masih berlaku dan apakah ada formulir yang lebih baru?
\end{verbatim}

\textbf{Email}

\begin{verbatim}
Subjek: [TGA] Konfirmasi kanal dan jadwal pendaftaran [tahap]: [Nama], [NIM]

Yth. Tim Administrasi Akademik DTAP,

Mohon konfirmasi tentang kanal, jadwal, dan formulir terbaru untuk pendaftaran [tahap]. Saya menemukan informasi di [sumber], yang saya akses pada [tanggal].

Jika informasi tersebut sudah berubah, mohon kirimkan tautan atau dokumen resmi terbaru.

Terima kasih atas bantuannya.

Hormat saya,
[Nama]
[NIM]
\end{verbatim}

\subsection{5.7 Menanyakan surat pengantar atau dokumen
administrasi}\label{menanyakan-surat-pengantar-atau-dokumen-administrasi}

\textbf{Langsung}

\begin{quote}
Saya ingin mengajukan {[}surat pengantar survei, surat penelitian, atau
dokumen lain{]}. Mohon informasi tentang persyaratan, format, kanal
pengajuan, dan tahapan pemrosesannya.
\end{quote}

\textbf{WhatsApp}

\begin{verbatim}
Mohon informasi tentang pengajuan [nama surat]. Saya ingin memastikan persyaratan, format berkas, kanal pengajuan, dan tahapan yang harus saya ikuti. Proposal saya [sudah atau belum] disetujui dosen pembimbing.
\end{verbatim}

\textbf{Email}

\begin{verbatim}
Subjek: [TGA] Permohonan informasi pengajuan [nama surat]: [Nama], [NIM]

Yth. Tim Administrasi Akademik DTAP,

Saya ingin mengajukan [nama surat] untuk keperluan [tujuan]. Proposal saya [sudah atau belum] disetujui dosen pembimbing.

Mohon informasi tentang dokumen persyaratan, format, kanal pengajuan, dan tahapan pemrosesan. Saya siap melampirkan [dokumen] jika diperlukan.

Terima kasih atas bantuannya.

Hormat saya,
[Nama]
[NIM]
\end{verbatim}

\subsection{5.8 Menanyakan pencatatan bimbingan, Progres, dan
prasidang}\label{menanyakan-pencatatan-bimbingan-progres-dan-prasidang}

\textbf{Langsung}

\begin{quote}
Saya ingin memastikan cara mencatat dan memverifikasi 10 kali bimbingan
serta nilai Progres 1 dan Progres 2 sebelum mendaftar prasidang. Sistem
apa yang digunakan, bukti apa yang perlu disiapkan, dan apakah Progres
otomatis dihitung sebagai bimbingan?
\end{quote}

\textbf{WhatsApp}

\begin{verbatim}
Mohon informasi, Bapak/Ibu. Untuk pendaftaran prasidang, bagaimana cara memverifikasi 10 kali bimbingan serta nilai Progres 1 dan Progres 2? Apakah Progres otomatis dihitung sebagai bimbingan atau perlu dicatat dan dikonfirmasi secara terpisah di SIMASTER?
\end{verbatim}

\textbf{Email}

\begin{verbatim}
Subjek: [TGA] Konfirmasi verifikasi bimbingan dan Progres sebelum prasidang: [Nama], [NIM]

Yth. Tim Administrasi Akademik DTAP,

Mohon konfirmasi tentang cara memverifikasi 10 kali bimbingan sebelum pendaftaran prasidang, termasuk sistem pencatatan dan bukti yang harus disiapkan.

Saya juga ingin memastikan apakah kegiatan Progres 1 dan Progres 2 otomatis dihitung sebagai bimbingan atau memerlukan pencatatan terpisah.

Terima kasih atas bantuannya.

Hormat saya,
[Nama]
[NIM]
\end{verbatim}

\subsection{5.9 Melaporkan kesalahan berkas atau meminta koreksi
prosedur}\label{melaporkan-kesalahan-berkas-atau-meminta-koreksi-prosedur}

\textbf{Langsung}

\begin{quote}
Saya menemukan {[}kesalahan atau status yang tidak sesuai{]} pada
{[}sistem atau berkas{]}. Saya sudah mencoba {[}langkah{]}. Apakah saya
perlu mengunggah ulang berkas, mengirimkannya melalui kanal lain, atau
menunggu perbaikan sistem?
\end{quote}

\textbf{WhatsApp}

\begin{verbatim}
Mohon bantuan, Bapak/Ibu. Saya menemukan [kesalahan] pada [sistem atau berkas] saat mengurus [urusan]. Saya sudah mencoba [langkah] pada [tanggal dan waktu].

Apakah saya perlu mengunggah ulang atau menggunakan kanal lain? Saya dapat mengirimkan tangkapan layar jika diperlukan.
\end{verbatim}

\textbf{Email}

\begin{verbatim}
Subjek: [TGA] Permohonan bantuan untuk [sistem atau berkas]: [Nama], [NIM]

Yth. Tim Administrasi Akademik DTAP,

Saat mengurus [urusan] pada [tanggal dan waktu], saya menemukan [uraian kesalahan]. Saya sudah mencoba [langkah].

Apakah saya perlu mengunggah ulang berkas, mengirimkannya melalui kanal lain, atau menunggu perbaikan sistem? Saya melampirkan tangkapan layar dan berkas terkait tanpa menyertakan data pribadi yang tidak diperlukan.

Terima kasih atas bantuannya.

Hormat saya,
[Nama]
[NIM]
\end{verbatim}

\subsection{5.10 Menanyakan yudisium, ETD, dan penyerahan
akhir}\label{menanyakan-yudisium-etd-dan-penyerahan-akhir}

\textbf{Langsung}

\begin{quote}
Saya ingin memastikan tahapan setelah sidang, termasuk yudisium, bebas
perpustakaan, unggah ETD, dan penyerahan dokumen akhir. Mohon informasi
tentang urutan, tenggat, dan dokumen yang harus saya siapkan.
\end{quote}

\textbf{WhatsApp}

\begin{verbatim}
Mohon informasi terbaru tentang tahapan setelah sidang, termasuk yudisium, bebas perpustakaan, unggah ETD, dan penyerahan dokumen akhir. Saya ingin memastikan urutan, tenggat, dan dokumen yang harus saya siapkan.
\end{verbatim}

\textbf{Email}

\begin{verbatim}
Subjek: [TGA] Konfirmasi tahapan pascasidang dan yudisium: [Nama], [NIM]

Yth. Tim Administrasi Akademik DTAP,

Mohon informasi terbaru tentang urutan tahapan setelah sidang, termasuk bebas perpustakaan, unggah ETD, formulir yudisium, penyerahan dokumen, dan tenggat setiap tahap.

Mohon informasi ini agar saya dapat menyiapkan berkas sesuai prosedur terbaru.

Terima kasih atas bantuannya.

Hormat saya,
[Nama]
[NIM]
\end{verbatim}

\section{6. Komunikasi dengan teman seangkatan, angkatan 23, dan pra-TA
angkatan
24}\label{komunikasi-dengan-teman-seangkatan-angkatan-23-dan-pra-ta-angkatan-24}

\subsection{6.1 Prinsip}\label{prinsip}

\begin{itemize}
\tightlist
\item
  Sampaikan pengalaman pribadi sebagai pengalaman, bukan sebagai aturan
  resmi.
\item
  Jika membagikan informasi administrasi, sertakan sumber dan tanggal
  akses.
\item
  Jangan menggurui angkatan 23 atau 24 hanya karena lebih dulu menjalani
  tahap tertentu.
\item
  Jangan meminta dokumen pribadi, isi bimbingan, atau data terbatas
  milik orang lain.
\item
  Bedakan bantuan akademik, bantuan administratif, dan keputusan yang
  hanya boleh dibuat dosen atau DTAP.
\item
  Hormati waktu teman. Ungkapan ``kalau sempat'' lebih baik daripada
  menuntut jawaban segera.
\end{itemize}

\subsection{6.2 Menanyakan pengalaman kepada angkatan 23 atau teman yang
sedang
skripsi}\label{menanyakan-pengalaman-kepada-angkatan-23-atau-teman-yang-sedang-skripsi}

\textbf{Bicara langsung}

\begin{quote}
Aku lagi menyiapkan {[}pra-TA atau TGA{]}. Kalau sempat, boleh cerita
bagaimana kamu mengurus {[}urusan spesifik{]}? Aturan resminya nanti
tetap aku cek ke DTAP.
\end{quote}

\textbf{WhatsApp}

\begin{verbatim}
Halo, [Nama]. Aku lagi menyiapkan [pra-TA atau TGA]. Kalau sempat, boleh cerita bagaimana kamu mengurus [urusan spesifik]?

Aturan resminya nanti tetap aku cek ke DTAP. Terima kasih, ya.
\end{verbatim}

\subsection{6.3 Meminta templat atau contoh
berkas}\label{meminta-templat-atau-contoh-berkas}

\textbf{Bicara langsung}

\begin{quote}
Kalau tidak keberatan, boleh minta contoh susunan berkas {[}nama
berkas{]}? Aku hanya ingin melihat formatnya, bukan menyalin isinya.
Data pribadi boleh disamarkan.
\end{quote}

\textbf{WhatsApp}

\begin{verbatim}
Kalau tidak keberatan, boleh minta contoh susunan berkas [nama berkas]? Aku hanya ingin melihat formatnya, bukan menyalin isinya. Data pribadi boleh disamarkan.
\end{verbatim}

\subsection{6.4 Menawarkan bantuan kepada pra-TA angkatan
24}\label{menawarkan-bantuan-kepada-pra-ta-angkatan-24}

\textbf{Bicara langsung}

\begin{quote}
Kalau kamu lagi menyiapkan pra-TA, aku bisa berbagi pengalaman tentang
cara membuat ringkasan topik, menghubungi dosen, atau mencatat arahan.
Untuk aturan resmi, tetap cek ke DTAP atau dosen karena informasinya
bisa berubah.
\end{quote}

\textbf{WhatsApp}

\begin{verbatim}
Halo, [Nama]. Kalau kamu lagi menyiapkan pra-TA, aku bisa berbagi pengalaman tentang ringkasan topik, cara menghubungi dosen, atau cara mencatat arahan.

Untuk aturan dan jadwal resmi, tetap cek ke DTAP atau dosen, ya. Pengalamanku belum tentu berlaku untuk angkatanmu.
\end{verbatim}

\subsection{6.5 Mengoreksi informasi yang mungkin
keliru}\label{mengoreksi-informasi-yang-mungkin-keliru}

\textbf{Bicara langsung}

\begin{quote}
Setahuku, informasi itu masih berdasarkan pengalaman {[}orang atau
sumber{]}, bukan pengumuman resmi. Aku menemukan {[}sumber dan tanggal
akses{]}. Sebaiknya kita cek lagi ke DTAP supaya tidak salah prosedur.
\end{quote}

\textbf{WhatsApp}

\begin{verbatim}
Setahuku, informasi itu masih berdasarkan pengalaman, bukan pengumuman resmi. Aku menemukan [sumber dan tanggal akses]. Sebaiknya kita cek lagi ke DTAP supaya tidak salah prosedur.
\end{verbatim}

\subsection{6.6 Mengatur koordinasi
teman}\label{mengatur-koordinasi-teman}

\textbf{Bicara langsung}

\begin{quote}
Supaya informasinya tidak simpang siur, bagaimana kalau sumber, tanggal
akses, dan statusnya kita catat di satu dokumen? Untuk keputusan
akademik, kita tetap mengikuti arahan dosen masing-masing.
\end{quote}

\textbf{WhatsApp}

\begin{verbatim}
Supaya informasinya tidak simpang siur, bagaimana kalau sumber, tanggal akses, dan statusnya kita catat di satu dokumen?

Untuk keputusan akademik, kita tetap mengikuti arahan dosen masing-masing, ya.
\end{verbatim}

\subsection{6.7 Menolak permintaan ketika sedang tidak
sempat}\label{menolak-permintaan-ketika-sedang-tidak-sempat}

\textbf{Bicara langsung}

\begin{quote}
Maaf, minggu ini aku lagi mengejar {[}pekerjaan{]}, jadi belum bisa
membantu sekarang. Kalau tidak mendesak, mungkin kita bisa bahas pada
{[}waktu{]}. Kalau mendesak, sebaiknya langsung tanyakan ke {[}kanal
resmi atau dosen{]}.
\end{quote}

\textbf{WhatsApp}

\begin{verbatim}
Maaf, minggu ini aku lagi mengejar [pekerjaan], jadi belum bisa membantu sekarang. Kalau tidak mendesak, mungkin kita bisa bahas pada [hari dan waktu]. Untuk aturan resmi, sebaiknya tetap tanyakan ke DTAP atau dosen.
\end{verbatim}

\subsection{6.8 Mengucapkan terima kasih kepada teman atau adik
tingkat}\label{mengucapkan-terima-kasih-kepada-teman-atau-adik-tingkat}

\textbf{Bicara langsung}

\begin{quote}
Terima kasih. Ceritamu membantu aku memahami prosesnya. Untuk bagian
resmi, aku tetap akan mengeceknya ke DTAP atau dosen.
\end{quote}

\textbf{WhatsApp}

\begin{verbatim}
Terima kasih banyak, [Nama]. Ceritamu membantu aku memahami prosesnya. Untuk bagian resmi, aku tetap akan mengeceknya ke DTAP atau dosen.
\end{verbatim}

\section{7. Aturan penjadwalan yang
etis}\label{aturan-penjadwalan-yang-etis}

\begin{itemize}
\tightlist
\item
  Utamakan hari dan jam kerja, misalnya pukul 08.00 sampai 17.00 waktu
  setempat, kecuali penerima sudah menetapkan kebiasaan lain.
\item
  Untuk dosen, tawarkan dua atau tiga pilihan waktu lengkap dengan hari,
  tanggal, jam, durasi, lokasi, dan agenda.
\item
  Jangan hanya bertanya, ``Bapak kapan ada waktu?'' Siapkan beberapa
  pilihan.
\item
  Konfirmasi ulang jadwal dan zona waktu jika pertemuan berlangsung
  daring.
\item
  Kirim bahan cukup awal. Jangan mengirim draf panjang beberapa menit
  sebelum pertemuan kecuali ada keadaan khusus.
\item
  Kirim satu pengingat sekitar sehari sebelumnya jika belum ada
  kebiasaan lain.
\item
  Jika tidak ada respons, beri jeda yang wajar sebelum mengirim satu
  pesan tindak lanjut. Jangan mengirim pesan beruntun atau bertanya,
  ``Sudah dibaca, Pak?''
\item
  Untuk administrasi, ikuti tenggat layanan resmi jika tersedia. Jika
  tidak tersedia, catat tanggal pengiriman dan kirim satu pesan tindak
  lanjut setelah jeda yang wajar.
\item
  Jika penerima mengubah jadwal, konfirmasi jadwal baru tanpa menyindir
  atau menunjukkan kejengkelan.
\item
  Jika mahasiswa yang mengubah jadwal, minta maaf secukupnya, berikan
  alasan singkat, dan tawarkan pengganti.
\item
  Jangan menjadwalkan bimbingan terlalu rapat. Beri waktu kepada dosen
  untuk membaca bahan dan kepada mahasiswa untuk mengerjakan tindak
  lanjut.
\item
  Jangan menganggap pesan WhatsApp sebagai persetujuan resmi untuk
  urusan yang harus masuk SIMASTER, formulir, surat, atau kanal
  akademik.
\item
  Setelah pertemuan, kirim ringkasan keputusan dan tindak lanjut pada
  hari yang sama atau sesegera mungkin.
\item
  Catat tanggal, agenda, keputusan, versi berkas, dan tindak lanjut
  dalam catatan komunikasi atau catatan bimbingan.
\end{itemize}

\section{8. Kiat praktis}\label{kiat-praktis}

\subsection{8.1 Sampaikan tujuan pesan dengan
jelas}\label{sampaikan-tujuan-pesan-dengan-jelas}

Jangan menulis ``Mohon dicek semuanya''. Sebutkan bagian yang perlu
dibaca dan keputusan yang Anda butuhkan:

\begin{verbatim}
Saya terutama memerlukan arahan tentang hal-hal berikut:
1. Apakah unit analisis tetap kecamatan?
2. Apakah baseline terbuka dapat digunakan untuk Progres 1?
3. Apakah asumsi estimasi dan bukti pemeriksaan independen untuk [modul] sudah memadai?
\end{verbatim}

\subsection{8.2 Gunakan nama versi
berkas}\label{gunakan-nama-versi-berkas}

Gunakan pola seperti \texttt{Nama\_Topik\_Bab3\_v03\_2026-08-28.pdf}.
Hindari nama seperti \texttt{final}, \texttt{final2}, atau
\texttt{terbaru\ banget} tanpa tanggal.

\subsection{8.3 Bedakan laporan, permintaan, dan
pertanyaan}\label{bedakan-laporan-permintaan-dan-pertanyaan}

\begin{itemize}
\tightlist
\item
  \textbf{Laporan:} apa yang sudah dikerjakan dan hasilnya.
\item
  \textbf{Permintaan:} keputusan atau tindakan yang Anda minta dari
  penerima.
\item
  \textbf{Pertanyaan:} hal yang belum diketahui dan perlu dijawab.
\end{itemize}

Satu pesan boleh memuat ketiganya, tetapi urutannya harus jelas.

\subsection{8.4 Gunakan ringkasan setelah
bimbingan}\label{gunakan-ringkasan-setelah-bimbingan}

\begin{verbatim}
Pak, berikut ringkasan bimbingan hari ini:

- Yang disepakati: [keputusan].
- Yang perlu direvisi: [bagian].
- Yang ditunda: [isu].
- Tindak lanjut saya: [pekerjaan], paling lambat [tanggal].

Jika ada catatan yang kurang tepat, mohon dikoreksi. Terima kasih, Pak.
\end{verbatim}

\subsection{8.5 Tetap sopan tanpa merendahkan
diri}\label{tetap-sopan-tanpa-merendahkan-diri}

Sopan tidak berarti menyebut diri ``hamba'', memakai sapaan ``Yang
Mulia'', meminta maaf berkali-kali, atau memberi pujian panjang. Rasa
hormat lebih terlihat dari identitas yang jelas, persiapan yang matang,
permintaan yang spesifik, dan tindak lanjut yang benar-benar dikerjakan.

\subsection{8.6 Jika dosen belum
membalas}\label{jika-dosen-belum-membalas}

Gunakan urutan berikut:

\begin{enumerate}
\def\labelenumi{\arabic{enumi}.}
\tightlist
\item
  Periksa apakah pesan berisi permintaan yang jelas.
\item
  Beri jeda yang wajar.
\item
  Kirim satu pesan tindak lanjut dengan tanggal pesan sebelumnya.
\item
  Tetap kerjakan baseline atau bahan lain yang tidak bergantung pada
  jawaban.
\item
  Jika menyangkut tenggat resmi, hubungi kanal akademik yang tepat.
  Untuk keputusan terkait isi penelitian, tetap libatkan dosen.
\end{enumerate}

\subsection{8.7 Jika arahan terasa berat atau mengubah
topik}\label{jika-arahan-terasa-berat-atau-mengubah-topik}

Jangan langsung menolak atau menyetujui tanpa memahami alasannya.
Tanyakan dampak perubahan tersebut terhadap pertanyaan, data, waktu, dan
klaim:

\begin{verbatim}
Pak, saya memahami arah revisinya. Agar saya bisa menyesuaikan jadwal, apakah perubahan ini juga mengubah [unit, data, atau metode], atau hanya cara pembahasannya?
\end{verbatim}

\subsection{8.8 Jika berkomunikasi tentang
data}\label{jika-berkomunikasi-tentang-data}

\begin{itemize}
\tightlist
\item
  Laporkan apa yang ditemukan, berhasil diekstrak, sudah diaudit, dan
  masih belum terverifikasi secara terpisah.
\item
  Jangan menulis ``data lengkap'' jika baru menemukan portal atau judul
  dataset.
\item
  Jangan menjanjikan hasil analisis sebelum unit, tahun, konstruk, dan
  mutu data diperiksa.
\item
  Bawa pilihan konstruksi indikator dan validasi berbasis sumber terbuka
  beserta konsekuensi masing-masing.
\end{itemize}

\subsection{8.9 Jika menerima informasi dari kakak tingkat atau
teman}\label{jika-menerima-informasi-dari-kakak-tingkat-atau-teman}

Catat tiga hal: siapa yang memberi informasi, kapan informasi diberikan,
dan apakah ada sumber resmi yang mendukungnya. Pengalaman angkatan 23
dapat membantu memperkirakan proses, tetapi tidak otomatis berlaku bagi
angkatan 22 atau 24.

\section{9. Bentuk pesan yang perlu
dihindari}\label{bentuk-pesan-yang-perlu-dihindari}

{\def\LTcaptype{none} % do not increment counter
\begin{longtable}[]{@{}
  >{\raggedright\arraybackslash}p{(\linewidth - 2\tabcolsep) * \real{0.5000}}
  >{\raggedright\arraybackslash}p{(\linewidth - 2\tabcolsep) * \real{0.5000}}@{}}
\toprule\noalign{}
\begin{minipage}[b]{\linewidth}\raggedright
Hindari
\end{minipage} & \begin{minipage}[b]{\linewidth}\raggedright
Ganti dengan
\end{minipage} \\
\midrule\noalign{}
\endhead
\bottomrule\noalign{}
\endlastfoot
``Pak, kapan bisa bimbingan?'' & Dua atau tiga pilihan tanggal, jam,
durasi, dan agenda \\
``Mohon dicek semuanya'' & Dua atau tiga bagian yang memerlukan
keputusan \\
``Sudah dibaca, Pak?'' & Tindak lanjut yang menyebut tanggal pengiriman
dan kebutuhan arahan \\
``Data ditolak'' ketika belum ada jawaban & ``Belum merespons sampai
{[}tanggal{]}'' \\
``Kata kakak tingkat pasti begini'' & ``Menurut pengalaman {[}sumber{]},
tetapi perlu diperiksa lagi ke DTAP'' \\
Mengirim draf tanpa versi atau konteks & Nama berkas, versi, tautan, dan
bagian yang perlu dibaca \\
Menyalahkan perubahan jadwal & Konfirmasi jadwal baru dengan tenang \\
Menghubungi semua kanal sekaligus & Pilih kanal yang sesuai dan tunggu
respons \\
Meminta bantuan tanpa batas kepada adik tingkat & Jelaskan bantuan yang
diminta dan beri ruang untuk menolak \\
\end{longtable}
}

\section{10. Catatan komunikasi}\label{catatan-komunikasi}

Gunakan tabel berikut secara ringkas. Jangan menyimpan informasi pribadi
yang tidak diperlukan.

{\def\LTcaptype{none} % do not increment counter
\begin{longtable}[]{@{}
  >{\raggedright\arraybackslash}p{(\linewidth - 12\tabcolsep) * \real{0.1429}}
  >{\raggedright\arraybackslash}p{(\linewidth - 12\tabcolsep) * \real{0.1429}}
  >{\raggedright\arraybackslash}p{(\linewidth - 12\tabcolsep) * \real{0.1429}}
  >{\raggedright\arraybackslash}p{(\linewidth - 12\tabcolsep) * \real{0.1429}}
  >{\raggedright\arraybackslash}p{(\linewidth - 12\tabcolsep) * \real{0.1429}}
  >{\raggedright\arraybackslash}p{(\linewidth - 12\tabcolsep) * \real{0.1429}}
  >{\raggedright\arraybackslash}p{(\linewidth - 12\tabcolsep) * \real{0.1429}}@{}}
\toprule\noalign{}
\begin{minipage}[b]{\linewidth}\raggedright
Tanggal
\end{minipage} & \begin{minipage}[b]{\linewidth}\raggedright
Penerima
\end{minipage} & \begin{minipage}[b]{\linewidth}\raggedright
Kanal
\end{minipage} & \begin{minipage}[b]{\linewidth}\raggedright
Tujuan
\end{minipage} & \begin{minipage}[b]{\linewidth}\raggedright
Berkas atau sumber
\end{minipage} & \begin{minipage}[b]{\linewidth}\raggedright
Respons
\end{minipage} & \begin{minipage}[b]{\linewidth}\raggedright
Keputusan atau tindak lanjut
\end{minipage} \\
\midrule\noalign{}
\endhead
\bottomrule\noalign{}
\endlastfoot
& & & & & & \\
& & & & & & \\
& & & & & & \\
\end{longtable}
}

\section{11. Batasan dan pembaruan}\label{batasan-dan-pembaruan}

\begin{itemize}
\tightlist
\item
  Sesuaikan panduan ini dengan pola komunikasi Bapak Retno yang terlihat
  dari interaksi langsung, bukan hanya dari dugaan berdasarkan publikasi
  beliau.
\item
  Jika beliau menetapkan kanal, waktu, format berkas, atau sapaan
  tertentu, ikuti kebiasaan tersebut.
\item
  Selalu periksa aturan DTAP, jadwal TGA, formulir, dan syarat
  pendaftaran melalui kanal resmi terbaru.
\item
  Panduan ini tidak menggantikan keputusan dosen pembimbing,
  administrasi akademik, atau kanal resmi TGA.
\end{itemize}

Chat

Assalamu'alaikum wr. wb.~ Selamat pagi, Pak Retno. Perkenalkan saya
Mahastri Laksita, mahasiswa PWK angkatan 22. Saya ingin meminta bapak
menjadi pembimbing Pra-TA untuk semester ini bapak, sebelumnya saya
meminta maaf tidak lulus Pra-TA tahun lalu dan tidak berkabar ke Bapak.~

Judul saya mengenai {[}insert{]} dan sudah membuat ringkasan satu
halaman yang berisi masalah, pertanyaan, metode awal, serta kebutuhan
data. Apakah saya boleh mengirimkannya dan meminta waktu untuk
berdiskusi tentang kemungkinan bimbingan? Saya berharap bisa bertemu
bapak di dtap dan menjelaskan Pra-TA saya secara tatap muka jika Pak
Retno berkenan.

Terima kasih, Pak

bubble chat pertama - intro - alasan tertarik mau dibimbing bapaknya

kedua - abstrak sejauh mana - ketertarikan saya - willing to explore
lagi jika diterima

\begin{center}\rule{0.5\linewidth}{0.5pt}\end{center}

\textbf{Peranmu:} Peneliti asisten yang \emph{evidence-based} di bidang
Urban Planning/Urban Design/Plannology (PWK).~~ \textbf{Tujuan:} Dari
sumber yang kuberikan, buat pemetaan menyeluruh tentang dosen pembimbing
(minat, lensa teori, metode, jaringan, gaya bimbing, dsb.)
\textbf{serta} turunkan \emph{riset-compass} (orientasi nilai yang
relevan untuk kerja akademik atau bahkan perkiraan preferensi politik
pribadi). Semua klaim harus ditopang bukti kutipan/tautan.

\textbf{Subjek \& Sumber Utama}~~ Nama:
\texttt{\{retno\ widodo\ dwi\ pramono,\ s.t.,\ m.sc.,\ ph.d.\}}~~ URL
inti ({[}https://archiplan.ugm.ac.id/id/dosen-program-pwk/{]},
{[}https://acadstaff.ugm.ac.id/Widodo{]},
{[}https://prisma.simaster.ugm.ac.id/Widodo\#course{]},
{[}https://prisma.simaster.ugm.ac.id/Widodo\#prisma{]},
{[}https://prisma.simaster.ugm.ac.id/Widodo\#prisma{]},
{[}https://prisma.simaster.ugm.ac.id/Widodo\#prisma{]},
{[}acadstaff.ugm.ac.id/Widodo\#{]},
{[}https://scholar.google.com/citations?user=MA-Ri3AAAAAJ\&hl=en{]},
{[}https://www.researchgate.net/profile/Retno-Pramono{]})~ ~~

Catatan: bisa ambil dari sumber di luar ini, jika dirasa butuh tambahan
info!

\begin{center}\rule{0.5\linewidth}{0.5pt}\end{center}

\section{Langkah Kerja (ikuti
berurutan)}\label{langkah-kerja-ikuti-berurutan}

\begin{enumerate}
\def\labelenumi{\arabic{enumi}.}
\item
  \textbf{Koleksi \& ringkas} 3--7 publikasi/aktivitas \textbf{terbaru}
  (\$\leq\$5 tahun) yang paling relevan: judul, tahun, outlet, peran
  (PI/co-author), \emph{topic}, \emph{theory}, \emph{method},
  \emph{data}, \emph{geografi/level skala} (tapak--kota--regional),
  \emph{temuan kunci}.
\item
  \textbf{Ekstraksi kata kunci} (top-15, dikelompokkan): \emph{tema}
  (TOD, perumahan, mobilitas aktif, NBS, heritage, resilience, smart
  city, informalitas, dsb.), \emph{kerangka teori} (Just City, Right to
  the City, New Urbanism, Landscape Urbanism, governance, dll.),
  \emph{metode} (GIS/spatial analysis, AHP/MCDA, Space Syntax, SEM/ML,
  ABM, PPGIS/FGD, etnografi, policy analysis, mixed).
\item
  \textbf{Peta metode \& data favorit} (frekuensi): desain
  (eksperimen/survei/studi kasus/etnografi/mixed), analitik
  (regresi/SEM/ML, network/space syntax), partisipatif (PPGIS, FGD,
  Delphi), sumber data (BPS, OSM, citra satelit, mobilitas, scraping,
  arsip kebijakan, survei lapangan).
\item
  \textbf{Lensa epistemologis \& kerangka}: positivis/empiris vs
  interpretif/kritis; teori yang konsisten dipakai/dikutip.
\item
  \textbf{Jaringan \& kolaborasi}: co-author berulang, lab/center, mitra
  pemerintah/NGO, akses data khas.
\item
  \textbf{Target outlet \& standar kualitas}: jurnal/konferensi sasaran,
  style (kuant rigour vs naratif kual), kedalaman analisis.
\item
  \textbf{Pendanaan \& agenda} (jika ada): skema hibah, tema prioritas.
\item
  \textbf{Gaya bimbing \& ekspektasi} Bisa anda perkirakan jika tidak
  ada~
\item ~
  \section{\texorpdfstring{\textbf{Keterbatasan praktis}: akses
  instrumen/data, ketersediaan waktu, periode
  sibuk.}{Keterbatasan praktis: akses instrumen/data, ketersediaan waktu, periode sibuk.}}\label{keterbatasan-praktis-akses-instrumendata-ketersediaan-waktu-periode-sibuk.}

  \section{Riset-Compass (bukan label politik
  pribadi)}\label{riset-compass-bukan-label-politik-pribadi}
\end{enumerate}

Untuk tiap sumbu di bawah, posisikan pada skala \textbf{-2 \ldots{} +2}
(bukti wajib), + \textbf{Confidence} (Low/Med/High):

\begin{itemize}
\tightlist
\item
  \textbf{Equity \(\leftrightarrow\) Efficiency} (keadilan/spatial
  justice vs efisiensi ekonomi)
\item
  \textbf{Participatory \(\leftrightarrow\) Technocratic}
  (PPGIS/FGD/CBPR vs expert-driven/teknokratik)
\item
  \textbf{State-led \(\leftrightarrow\) Market-led} (preferensi
  kebijakan/pendekatan implementasi)
\item
  \textbf{Conservation/Heritage \(\leftrightarrow\) Redevelopment}
\item
  \textbf{Nature-based \(\leftrightarrow\) Grey infrastructure}
\item
  \textbf{Interpretive/Critical \(\leftrightarrow\)
  Positivist/Empirical}
\item
  \textbf{Low-tech/Community \(\leftrightarrow\) Smart-city/High-tech}
\item
  \textbf{Epistemik} (konstruktivis/interpretif \(\leftrightarrow\)
  positivis/empiris)
\item
  \textbf{Sikap terhadap risiko \& ketidakpastian} (berani eksplor
  \(\leftrightarrow\) konservatif metodologis)~
\end{itemize}

\begin{quote}
Untuk setiap sumbu: 2--3 bullet penjelasan kenapa skornya demikian.
\end{quote}

\begin{center}\rule{0.5\linewidth}{0.5pt}\end{center}

\section{Keluaran yang Harus Dibuat}\label{keluaran-yang-harus-dibuat}

\begin{enumerate}
\def\labelenumi{\arabic{enumi}.}
\tightlist
\item
  \textbf{Executive Summary} (\$\leq\$120 kata) --- gambaran minat,
  metode andalan, dan 2--3 implikasi untuk strategi judul skripsi.
\item
  \textbf{Advisor Map (9 komponen)} --- dalam tabel ringkas.
\item
  \textbf{Riset-Compass} --- daftar sumbu + skor
\item
  \textbf{Judul Skripsi yang ``Match'' (3 opsi)} --- format: ~ ~ -
  \emph{{[}Fenomena/Objek{]} + {[}Konteks{]} + {[}Kerangka{]} +
  {[}Metode/Data{]} + {[}Tujuan/kontribusi{]}} ~ ~ - Sertakan
  \textbf{alasan kecocokan} (2--3 kalimat) \& \textbf{tingkat risiko}
  (waktu/data).
\item
  \textbf{Fit Matrix (skor 1--5)} untuk tiap judul: Kesesuaian topik
  \textbar{} metode \textbar{} ketersediaan data \textbar{} nilai
  kontribusi \textbar{} risiko timeline.
\end{enumerate}

\textbf{Format Output} - Gunakan heading jelas, tabel untuk Advisor Map
\& Fit Matrix, bullet untuk bukti. - Cantumkan tautan langsung (URL)
pada setiap referensi/bukti. - Jika bukti tidak memadai untuk suatu
sumbu/komponen, tulis: \emph{``Belum terverifikasi (butuh sumber
tambahan: \ldots)''}.

\textbf{Gaya \& Kualitas} - Ringkas, spesifik, tidak spekulatif. -
Bahasa Indonesia, istilah teknis PWK OK.

\end{document}
